\documentclass[11pt]{article}
\usepackage[margin=1in]{geometry}
\usepackage{amsmath,amssymb,amsthm}
\usepackage[hidelinks]{hyperref}
\newtheorem{theorem}{Theorem}
\newcommand{\C}{\mathbb C}
\newcommand{\A}{\mathcal A}
\newcommand{\norm}[1]{\left\lVert#1\right\rVert}
\title{An injective C*-power series violating the Bieberbach coefficient bound}
\author{Research proof and Lean formalization}
\date{October 27, 2026}
\begin{document}
\maketitle
\begin{abstract}
We give a normalized power series on $C([0,1],\C)$ which converges in norm
at every element and is globally injective, but whose second coefficient
violates the proposed C*-algebraic Bieberbach order bound. Its formal
derivative series also converges everywhere. Injectivity follows from a
continuous exponential period anchored at zero; injectivity of individual
scalar fibers is not assumed.
\end{abstract}

Krishna's Conjecture~1.2 in \cite{Krishna} proposes the following extension
of the scalar Bieberbach theorem. If $\A$ is a unital C*-algebra and
$F(u)=u+\sum_{n\geq2}a_nu^n$ converges in norm and is injective on
$\{u:\norm u<1\}$, then
\begin{equation}\label{bound}
 a_na_n^*\leq n^2 1,\qquad a_n^*a_n\leq n^2 1\qquad(n\geq2).
\end{equation}
Conjecture~1.4 adds norm convergence of the derivative series. The next
example violates both general-coefficient assertions. The separate
invertible-coefficient variants are not claimed refuted.

\begin{theorem}\label{main}
Let $\A=C([0,1],\C)$ with its supremum norm, and let $a(t)=8t$. Put
\begin{equation}\label{series}
 c_n=\frac{a^{n-1}}{n!}\quad(n\geq1),\qquad
 F(u)=\sum_{n\geq1}c_nu^n.
\end{equation}
This series and its formal derivative series converge in the norm of $\A$
for every $u\in\A$. The map $F:\A\to\A$ is injective. Nevertheless
$c_2(1)=4$, so both inequalities in \eqref{bound} fail at $n=2$.
\end{theorem}

\begin{proof}
For every $u\in\A$,
\[
 \sum_{n\geq1}\norm{c_nu^n}
 \leq\norm u\sum_{k\geq0}\frac{\norm{au}^k}{(k+1)!}<\infty.
\]
The derivative series is dominated by an exponential series as well.
Multiplication of \eqref{series} by $a$ and reindexing the absolutely
convergent exponential series gives the identity in $\A$
\begin{equation}\label{identity}
 aF(u)=\exp(au)-1.
\end{equation}
Here the exponential is the Banach-algebra exponential, whose scalar
evaluation is the usual complex exponential. Evaluation of the norm-convergent
series at $t=0$ also gives $F(u)(0)=u(0)$. In particular,
\[
 F(u)(t)=\frac{\exp(8tu(t))-1}{8t}\quad(t>0).
\]

Suppose $F(u)=F(v)$. By \eqref{identity}, the continuous function
$h(t)=8t(u(t)-v(t))$ satisfies $\exp(h(t))=1$ for every $t$, and $h(0)=0$.
Thus its image lies in the discrete set $2\pi i\mathbb Z$.
Connectedness of $[0,1]$ forces $h=0$. Since $8t\ne0$ for $t>0$, we obtain
$u(t)=v(t)$ there, and the equality at $0$ follows from $F(u)(0)=F(v)(0)$.
Hence $u=v$.

Finally $c_1=1$ and $c_2(t)=4t$. Evaluation of either claimed order bound
at $t=1$ would give $16\leq4$, a contradiction.
\end{proof}

\paragraph{Why the scalar theorem does not apply.}
The source requires injectivity of a map on the unit ball of the algebra.
It does not require that every scalar fiber be univalent. A collision of two
continuous inputs in this example would produce a continuous integer-valued
exponential period, which must vanish because it vanishes at the endpoint.
Some individual scalar fibers are noninjective. The classical scalar
Bieberbach theorem is therefore not contradicted. The coefficients $c_n$
for $n\geq2$ vanish at $0$ and are not invertible elements of $\A$; the
source's invertible-coefficient restrictions cannot be silently discarded.

\paragraph{Formal scope and review status.}
The Lean certificate has 28 declarations in four modules. It proves actual
Banach-norm summability, the exponential identity, scalar evaluation, global
injectivity, the literal coefficient-order violation, and closed negations
of the two source assertions specialized to the genuine commutative unital
C*-algebra $C([0,1],\C)$. No convergence, injectivity or fiber-choice
hypothesis remains. The endpoints are
\[
\texttt{source\_counterexample},\qquad
\texttt{source\_strong\_counterexample}
\]
in namespace \texttt{CStarBieberbach}. The isolated strict replay uses
Lean~4.34.1 and mathlib revision
\[\texttt{d13f23b723b8a846827a245b89c10fc7d3f11612},\]
with only \texttt{propext}, \texttt{Classical.choice} and \texttt{Quot.sound}.
It checks the fixed coefficient $8t$; no separately formalized theorem about
all parameter values is claimed.

Independent preliminary source, mathematical, bounded prior-art and
substantive-eligibility review passes; complete final review is pending.
The source was read in the author-deposited Zenodo version
(doi:10.5281/zenodo.6917925), whose text agrees with the source reading of
Preprints version~1. Direct Preprints PDF retrieval remains unavailable.
Discovery credit remains zero. The two variants use one
construction and constitute at most one prospective outcome. Classical
exponential-series and connectedness facts and helper declarations are not
additional discoveries. No worldwide priority claim is made.

\begin{thebibliography}{1}
\bibitem{Krishna} K.~Mahesh Krishna,
\emph{C*-algebraic Bieberbach, Robertson, Lebedev--Milin, Zalcman, Krzyz and
Corona Conjectures}, Preprints (2022), version~1,
Conjectures~1.2 and~1.4, pp.~2--3.
\href{https://doi.org/10.20944/preprints202208.0140.v1}{doi:10.20944/preprints202208.0140.v1}.
\end{thebibliography}
\end{document}
